\begin{circuitikz}[
    american,
    voltage dir=RP,
    >=latex,
    %>={Latex[length=5mm, width=2mm]},
    alt={}
    ]

    % Define the number of rungs in the diagram
    \def\numRungs{1}

    % Calculate the length of the vertical power rails dynamically based on number of rungs
    \pgfmathsetmacro{\railLength}{-2 * \numRungs}

    %======================================================================================================
    % Component Grid
    % Spacing = 3x, 2y
    %======================================================================================================
    \foreach \i in {0,...,5} {
        \foreach \j in {0,...,20} { % Change to 0,...,19 if you need 20 individual Y points down to -40
        
        % Calculate spatial positions
        \pgfmathsetmacro{\px}{\i * 3}
        \pgfmathsetmacro{\py}{-1 - (\j * 2)}
        
        % Name coordinate using loop indices: (C-column-row) e.g., (C-0-0), (C-5-10)
        \coordinate (C-\i-\j) at (\px, \py);
        
        % OPTIONAL: Uncomment the 2 lines below to display grid labels while building
        % \fill[red] (C-\i-\j) circle (1.5pt);
        %\node[anchor=south west, scale=0.5, gray] at (C-\i-\j) {(\i,\j)};
        }
    }

    \begin{pgfonlayer}{ft}
        
        \draw[
            BakerBlack,
            very thick,
            name path=L1
        ]
        (0,0)
        -- (0,\railLength);

        \draw[
            BakerBlack,
            very thick,
            name path=L2
        ]
        (15,0)
        -- (15,\railLength);

    \end{pgfonlayer}

    \begin{pgfonlayer}{main}

        %=================================================================
        % Rung 1
        %=================================================================
        \draw[
            BakerBlack
        ]
        (C-0-0)
        to[C, color=BakerBlack, l={{\small Auto\_Mode}}] (C-1-0)
        -- (C-3-0)
        pic (T1) {TON};

        \draw[
            BakerBlack
        ]
        (T1-OUT)
        -- (C-5-0);

        \node[text=BakerBlack, anchor=south] at (T1-TOP) {{\small RTO}};
        \node[text=BakerBlack, anchor=north] at (T1-TOP) {{\small AutoTimer}};
        \node[text=BakerBlack, anchor=west] at (T1-PRE) {{\small 28800}};
        \node[text=BakerBlack, anchor=west] at (T1-ACC) {{\small 15608}};

    \end{pgfonlayer}

    %=================================================================
    % Background layer - highlight any logic that should be true
    % UpnorthGreen, opacity=0.4, line width=10pt
    %=================================================================
    \begin{pgfonlayer}{bg}
        
        %=================================================================
        % Rung 1
        %=================================================================
        % \draw[
        %     UpNorthGreen,
        %     opacity=0.4,
        %     line width=10pt,
        % ]
        % (C-0-0)
        % -- (C-1-0);

        % \draw[
        %     UpNorthGreen,
        %     opacity=0.4,
        %     line width=10pt,
        % ]
        % (T1-EN1)
        % -- (T1-EN2);

        % \draw[
        %     UpNorthGreen,
        %     opacity=0.4,
        %     line width=10pt,
        % ]
        % (T1-DN1)
        % -- (T1-DN2);

    \end{pgfonlayer}
    
\end{circuitikz}